% Unit 100 terjemahan Bahasa Indonesia dari chapter7.tex baris 2141--2305.
% Sumber: Methods of Algebra, Volume 2, commit
% 9a5803ff77dd3257484cb177f851a73770a59dd3 (CC BY 4.0).
% stable-unit-id: o014.aljabr2.chapter7.module-descent
% Cakupan: Bagian 7.9 lengkap, ``Penerapan: Penurunan Modul'';
% berhenti tepat sebelum Bagian 7.10 pada baris 2306.

% segment-id: o014.li2.chapter7.module-descent.h001
\section{Penerapan: Penurunan Modul}\label{sec:descent-modules}
\hypertarget{o014.aljabr2.chapter7.module-descent}{ }
% segment-id: o014.li2.chapter7.module-descent.p001
Dalam geometri, untuk pasangan adjoin yang diberikan
% segment-id: o014.li2.chapter7.module-descent.q001
$\begin{tikzcd}
	F: \mathcal{C} \arrow[shift left, r] & \mathcal{D}: G \arrow[shift left, l]
\end{tikzcd}$,
proses merekonstruksi informasi pada $\mathcal{C}$ dari komodul pada
$\mathcal{D}$ di bawah koaksi komonad yang bersesuaian biasanya disebut
\emph{penurunan}. Bagian ini menjelaskannya dalam teori modul.
\index{desen@desen (descent)}

% segment-id: o014.li2.chapter7.module-descent.p002
Tinjau homomorfisme gelanggang komutatif $K\to L$. Funktor pelupa yang
bersesuaian $\mathcal{F}_{L|K}:L\dcate{Mod}\to K\dcate{Mod}$ dapat
ditempatkan dalam pasangan adjoin
% segment-id: o014.li2.chapter7.module-descent.q002
\[\begin{tikzcd}
	L \dotimes{K} (\cdot): K\dcate{Mod} \arrow[shift left, r] & L\dcate{Mod} : \mathcal{F}_{L|K}, \arrow[shift left, l]
\end{tikzcd}\]
Kita akan menerapkan hasil pada paruh kedua \S\ref{sec:Morita} untuk meninjau
komonad yang bersesuaian pada $L\dcate{Mod}$. Lebih tepatnya, ambil
$\Bbbk=K$ dan terapkan ``perubahan gelanggang'' dari
Contoh~\ref{eg:comonad-change-ring}. Selain perbedaan notasi, di sini kita
memakai modul kiri, bukan modul kanan. Karena hanya gelanggang komutatif yang
dipertimbangkan, hal ini tidak membawa perbedaan substantif.

% segment-id: o014.li2.chapter7.module-descent.p003
Berdasarkan deskripsi dalam Contoh~\ref{eg:comonad-change-ring}, komonad yang
ditentukan pasangan adjoin tersebut pada $L\dcate{Mod}$ ditulis
$(\mathbf{L},\delta,\epsilon)$ dalam bagian ini, dengan
% segment-id: o014.li2.chapter7.module-descent.q003
\[ \mathbf{L} = L \dotimes{K} \mathcal{F}_{L|K}(\cdot), \quad \delta: \mathbf{L} \to \mathbf{L}^2, \quad \epsilon: \mathbf{L} \to \identity. \]
Selanjutnya, sesuai kebiasaan, funktor pelupa $\mathcal{F}_{L|K}$ akan
dihilangkan dari berbagai rumus untuk menyederhanakan notasi.

% segment-id: o014.li2.chapter7.module-descent.p004
Komodul bagi komonad $\mathbf{L}$ adalah data $(M,a)$, dengan $M$ modul-$L$
dan $a:M\to L\dotimes{K}M$ homomorfisme modul-$L$, yang disyaratkan membuat
diagram berikut komutatif:
% segment-id: o014.li2.chapter7.module-descent.q004
\[\begin{tikzcd}
	\ell \otimes 1 \otimes m \arrow[phantom, r, "\in" description] & L \dotimes{K} L \dotimes{K} M & L \dotimes{K} M \arrow[l, "{\mathbf{L}a = \identity_L \otimes a}"' inner sep=0.6em] \\
	\ell \otimes m \arrow[phantom, r, "\in" description] \arrow[mapsto, u] & L \dotimes{K} M \arrow[u, "\delta_M"] & M \arrow[l, "a"] \arrow[u, "a"']
\end{tikzcd} \quad \begin{tikzcd}
	\ell m \arrow[phantom, d, "\in" description, sloped] & \ell \otimes m \arrow[phantom, d, "\in" description, sloped] \arrow[mapsto, l] \\
	M & L \dotimes{K} M \arrow[l, "\epsilon_M"'] \\
	& M. \arrow[lu, "\identity_M"] \arrow[u, "a"']
\end{tikzcd}\]

% segment-id: o014.li2.chapter7.module-descent.p005
Semua komodul $(M,a)$ di bawah komonad $\mathbf{L}$ membentuk kategori
$L\dcate{Mod}^{\mathbf{L}}$. Versi dual Lema~\ref{prop:T-comparison}
memberikan faktorisasi
% segment-id: o014.li2.chapter7.module-descent.q005
\[ L \dotimes{K} (\cdot) = \left[ K\dcate{Mod} \xrightarrow{\mathbb{K}} L\dcate{Mod}^{\mathbf{L}} \xrightarrow{U^{\mathbf{L}}} L\dcate{Mod} \right] \; \text{sebagai suatu komposisi}. \]
Funktor $U^{\mathbf{L}}$ adalah funktor pelupa $(M,a)\mapsto M$, sedangkan
funktor $\mathbb{K}$ memetakan objek $N$ ke $(L\dotimes{K}N,a)$ dengan
% segment-id: o014.li2.chapter7.module-descent.q006
\begin{equation}\label{eqn:N-comonad-action}
	\begin{aligned}
		a: L \dotimes{K} N & \to L \dotimes{K} (L \dotimes{K} N) \quad \text{(homomorfisme modul-$L$)} \\
		\ell \otimes x & \mapsto \ell \otimes 1 \otimes x.
	\end{aligned}
\end{equation}

% segment-id: o014.li2.chapter7.module-descent.p006
Pertanyaan utamanya ialah apakah pasangan adjoin tersebut komonadik; dengan
kata lain, apakah $\mathbb{K}$ merupakan ekuivalensi?

% segment-id: o014.li2.chapter7.module-descent.d001
\begin{definition}
	\index{datar setia@datar setia (faithfully flat)}
	Untuk homomorfisme gelanggang komutatif $K\to L$, jika
	$L\dotimes{K}(\mathord\cdot):K\dcate{Mod}\to L\dcate{Mod}$ merupakan
	funktor eksak setia, maka $L$, sebagai aljabar-$K$ komutatif, disebut
	\emph{datar setia}.
\end{definition}

% segment-id: o014.li2.chapter7.module-descent.th001
\begin{theorem}[Penurunan datar]\label{prop:flat-descent}
	\index{desen!datar@datar (flat)}
	Misalkan $L$ aljabar-$K$ komutatif yang datar setia. Maka pasangan adjoin
	$\left(L\dotimes{K}(\mathord\cdot),\mathcal{F}_{L|K}\right)$ bersifat
	komonadik. Dalam hal ini, suatu kuasi-invers $\mathbb{L}$ dari
	$K\dcate{Mod}\xrightarrow{\mathbb{K}}L\dcate{Mod}^{\mathbf{L}}$ dapat
	dideskripsikan melalui diagram penyama dalam $K\dcate{Mod}$:
	% segment-id: o014.li2.chapter7.module-descent.q007
	\[\begin{tikzcd}[row sep=tiny]
		\mathbb{L}(M, a) \arrow[r] & M \arrow[shift left, r, "a"] \arrow[shift right, r, "{m \mapsto 1 \otimes m}"'] & L \dotimes{K} M ,
	\end{tikzcd} \quad (M, a) \in \Obj\left(L\dcate{Mod}^{\mathbf{L}}\right). \]
\end{theorem}
% segment-id: o014.li2.chapter7.module-descent.pf001
\begin{proof}
	Proposisi~\ref{prop:Morita-descent} memastikan sifat komonadik, sedangkan
	versi dual Teorema~\ref{prop:Beck} telah memuat deskripsi funktor kuasi-
	inversnya.
\end{proof}

% segment-id: o014.li2.chapter7.module-descent.r001
\begin{remark}
	Perlu ditekankan bahwa pada tingkat kategori turunan
	$\cate{D}(K\dcate{Mod})$ dan $\cate{D}(L\dcate{Mod})$, meskipun pasangan
	adjoin yang dibahas juga mempunyai versi turunan seperti dalam
	Contoh~\ref{eg:change-of-ring-adjunction} dan
	Contoh~\ref{eg:change-of-rings}, tidak ada teorema penurunan yang
	bersesuaian.
\end{remark}

% segment-id: o014.li2.chapter7.module-descent.p007
Bukti penurunan datar ternyata sangat sederhana. Namun, dalam praktik,
komonad atau data penurunan perlu dideskripsikan dalam bentuk yang mudah
digunakan; hal inilah yang menjadi tema \S\ref{sec:descent-Galois}. Pada sisa
bagian ini, definisi $L\dcate{Mod}^{\mathbf{L}}$ akan disusun ulang dalam
bentuk lain agar mudah dibandingkan dengan pustaka lain. Pada dasarnya,
ini adalah penulisan konkret ulang dari
Proposisi~\ref{prop:Morita-comonad} (dalam versi modul kiri, dengan
$P=L=P^\vee$), meskipun prosesnya agak berliku.

% segment-id: o014.li2.chapter7.module-descent.p008
Untuk $i=1,2$, tuliskan $\iota_i:L\to L\dotimes{K}L$ bagi homomorfisme yang
memasukkan unsur ke slot tensor ke-$i$. Untuk sembarang modul-$L$ $M$,
terdapat isomorfisme kanonik modul-$L$
% segment-id: o014.li2.chapter7.module-descent.q008
\begin{align*}
	L \dotimes{K} M & \rightiso (L \dotimes{K} L) \dotimes{L, \iota_2} M \\
	\ell \otimes m & \mapsto \ell \otimes 1 \otimes m,
\end{align*}
dengan ruas kanan menjadi modul-$L$ melalui $\iota_1$. Menurut adjungsi
standar dalam teori modul, memberikan homomorfisme modul-$L$
$a:M\to L\dotimes{K}M$ sama artinya dengan memberikan homomorfisme modul-
$L\dotimes{K}L$
% segment-id: o014.li2.chapter7.module-descent.q009
\begin{equation}\label{eqn:flat-descent-map-aux}\begin{aligned}
		\tilde{a}: (L \dotimes{K} L) \dotimes{L, \iota_1} M & \to (L \dotimes{K} L) \dotimes{L, \iota_2} M \\
		1 \otimes 1 \otimes m & \mapsto \sum_{i=1}^n \ell_i \otimes 1 \otimes m_i ;
\end{aligned}\end{equation}
di sini ditulis $a(m)=\sum_{i=1}^n\ell_i\otimes m_i$. Karena
$\iota_1|_K=\iota_2|_K$, bila $M$ berasal dari modul-$K$ $N$, substitusi
cermat dari \eqref{eqn:N-comonad-action} menunjukkan bahwa $\tilde a$
menjadi
$\identity:L\dotimes{K}L\dotimes{K}N\to L\dotimes{K}L\dotimes{K}N$.

% segment-id: o014.li2.chapter7.module-descent.p009
Selanjutnya, untuk menyingkat notasi, tuliskan
$L^{\otimes n}:=L\dotimes{K}\cdots\dotimes{K}L$ (sebanyak $n$ faktor).
Lupakan untuk sementara $L\dcate{Mod}^{\mathbf{L}}$ dan tinjau sembarang
$\tilde a:L^{\otimes2}\dotimes{L,\iota_1}M\to
L^{\otimes2}\dotimes{L,\iota_2}M$. Menurut sifat universal hasil kali tensor
\cite[Proposisi 7.3.4]{Li1}, memberikan sepasang homomorfisme aljabar-$K$
komutatif $f,g:L\to E$ sama artinya dengan memberikan homomorfisme
$f\otimes g:L\dotimes{K}L\to E$. Karena itu, dari $\tilde a$ dapat
didefinisikan homomorfisme modul-$E$
% segment-id: o014.li2.chapter7.module-descent.q010
\begin{equation*}
	\tilde{a}_{gf} := E \dotimes{f \otimes g} \tilde{a}: E \dotimes{L, f} M \to E \dotimes{L, g} M.
\end{equation*}

% segment-id: o014.li2.chapter7.module-descent.p010
Ketika $(E,f,g)$ berubah, kita ingin menganalisis kompatibilitas antara
homomorfisme tersebut. Rumus di atas menunjukkan bahwa
$(E,f,g):=(L\dotimes{K}L,\iota_1,\iota_2)$ bersifat ``universal'' untuk
persoalan ini dan bersesuaian dengan $\tilde a$, sedangkan $(E,f,g)$ umum
diperoleh dengan menerapkan $E\dotimes{f\otimes g}(\mathord\cdot)$ padanya.
Untuk menguraikan kasus ``universal'', bagi $1\leq i\neq j\leq3$ tuliskan:
% segment-id: o014.li2.chapter7.module-descent.q011
\begin{center}\begin{tabular}{|c|c|} \hline
		$\iota^3_{ij}: L^{\otimes 2} \to L^{\otimes 3}$ & penyisipan ke slot tensor $(i,j)$ \\
		$\iota^3_i: L \to L^{\otimes 3}$ & penyisipan ke slot tensor ke-$i$ \\ \hline
\end{tabular}\end{center}

% segment-id: o014.li2.chapter7.module-descent.p011
Dari \eqref{eqn:flat-descent-map-aux}, definisikan homomorfisme modul-
$L^{\otimes3}$
% segment-id: o014.li2.chapter7.module-descent.q012
\[ \tilde{a}_{ji} := L^{\otimes 3} \dotimes{L, \iota^3_{ij}}\tilde{a} : \; L^{\otimes 3} \dotimes{L, \iota^3_i} M \to L^{\otimes 3} \dotimes{L, \iota^3_j} M. \]
Di sini digunakan $\iota^3_{ij}\iota_1=\iota^3_i$ dan
$\iota^3_{ij}\iota_2=\iota^3_j$. Bila $M$ berasal dari modul-$K$, kita dapat
mengidentifikasi
% segment-id: o014.li2.chapter7.module-descent.q013
\[ \tilde{a} \;\text{dengan}\; \identity_{L^{\otimes 2} \dotimes{K} N}, \quad \tilde{a}_{ji} \;\text{dengan}\; \identity_{L^{\otimes 3} \dotimes{K} N}. \]
Dalam keadaan ini, jelas bahwa
$\tilde a_{31}=\tilde a_{32}\tilde a_{21}$. Inilah syarat kompatibilitas yang
dicari. Kegunaannya yang penting ialah
% segment-id: o014.li2.chapter7.module-descent.q014
\begin{equation}\label{eqn:cocycle-condition-descent}
	\tilde{a}_{31} = \tilde{a}_{32} \tilde{a}_{21} \implies
	\begin{array}{l}
		\forall f, g, h: L \xrightarrow{K\text{-aljabar}} E, \\
		\tilde{a}_{hf} = \tilde{a}_{hg} \tilde{a}_{gf};
	\end{array}
\end{equation}
dengan ruas kiri merupakan ``versi universal'' dari ruas kanan.

% segment-id: o014.li2.chapter7.module-descent.d002
\begin{definition}[Data penurunan]
	\index{datum desen@datum desen (descent datum)}
	Definisikan $\mathrm{Desc}_{K\to L}$ sebagai kategori berikut. Objeknya
	adalah data $(M,\tilde a)$, dengan $M$ modul-$L$ dan
	$\tilde a:L^{\otimes2}\dotimes{L,\iota_1}M\to
	L^{\otimes2}\dotimes{L,\iota_2}M$ memenuhi
	\begin{description}
		\item[Syarat isomorfisme] $\tilde a$ adalah isomorfisme modul-
		$L^{\otimes2}$;
		\item[Syarat kosikel] $\tilde a_{31}=\tilde a_{32}\tilde a_{21}$.
	\end{description}

	Morfisme dari $(M,\tilde a)$ ke $(M',\tilde a')$ adalah homomorfisme
	modul-$L$ $f:M\to M'$ yang memenuhi
	$(\identity^{\otimes2}\otimes f)\tilde a=
	\tilde a'(\identity^{\otimes2}\otimes f)$. Data $(M,\tilde a)$ ini
	disebut \emph{data penurunan}.
\end{definition}

% segment-id: o014.li2.chapter7.module-descent.p012
Pembahasan sebelum \eqref{eqn:cocycle-condition-descent} menunjukkan bahwa
$L\dotimes{K}(\mathord\cdot)$ memberikan funktor
$K\dcate{Mod}\to\mathrm{Desc}_{K\to L}$. Sekarang kita hubungkan data
penurunan dengan $L\dcate{Mod}^{\mathbf{L}}$.

% segment-id: o014.li2.chapter7.module-descent.le001
\begin{lemma}\label{prop:descent-data-vs-comodule}
	Terdapat isomorfisme kategori dari $L\dcate{Mod}^{\mathbf{L}}$ ke
	$\mathrm{Desc}_{K\to L}$ yang pada objek adalah
	$(M,a)\leftrightarrow(M,\tilde a)$ dan pada morfisme didefinisikan dengan
	cara baku.

	Selain itu, syarat isomorfisme dalam definisi
	$\mathrm{Desc}_{K\to L}$ dapat diganti dengan syarat berikut: definisikan
	homomorfisme $m:L^{\otimes2}\to L$ melalui
	$m(\ell\otimes\ell')=\ell\ell'$. Endomorfisme $M\to M$ yang diberikan oleh
	$L\dotimes{L^{\otimes2},m}\tilde a$ adalah $\identity_M$.
\end{lemma}
% segment-id: o014.li2.chapter7.module-descent.pf002
\begin{proof}
	Diketahui bahwa memberikan $\tilde a$ sama artinya dengan memberikan
	$a:M\to L\dotimes{K}M$. Persoalannya ialah mencocokkan syarat pada data
	penurunan $\tilde a$ dengan syarat pada $a$ agar $(M,a)$ menjadi komodul.
	Di sinilah bagian kedua pernyataan akan digunakan.

	Syarat komodul adalah hukum koasosiatif dan hukum kounit. Untuk hukum
	pertama, ekuivalensi antara
	$\tilde a_{31}=\tilde a_{32}\tilde a_{21}$ dan hukum koasosiatif merupakan
	pemeriksaan mekanis yang diserahkan kepada pembaca. Untuk hukum kedua,
	tuliskan $\varphi$ bagi komposisi
	$M\xrightarrow{a}L\dotimes{K}M\xrightarrow{\epsilon_M}M$. Hukum kounit
	sama artinya dengan $\varphi=\identity$. Dengan mengingat bahwa
	$\epsilon_M(\ell\otimes m)=\ell m$, pemeriksaan sederhana serupa memberikan
	% segment-id: o014.li2.chapter7.module-descent.q015
	\[ \varphi = L \dotimes{L^{\otimes 2}, m} \tilde{a}. \]
	Selanjutnya, kita tunjukkan bahwa dengan asumsi
	$\tilde a_{31}=\tilde a_{32}\tilde a_{21}$, $\tilde a$ adalah isomorfisme
	jika dan hanya jika $\varphi=\identity_M$. Hal ini menyelesaikan seluruh
	bukti.

	Misalkan $\tilde a$ isomorfisme. Maka $\varphi$ juga isomorfisme. Dalam
	\eqref{eqn:cocycle-condition-descent}, ambil $E=L$ dan
	$f=g=h=\identity_L$. Mudah diperiksa bahwa ruas kanannya menjadi
	$\varphi=\varphi^2$, sehingga $\varphi=\identity_M$.

Sebaliknya, misalkan $\varphi=\identity_M$. Untuk sembarang homomorfisme
aljabar-$K$ $f:L\to E$, lakukan perubahan skalar pada $\tilde a$ sepanjang kedua
	lintasan diagram komutatif
	% segment-id: o014.li2.chapter7.module-descent.q016
	\[\begin{tikzcd}
		L^{\otimes 2} \arrow[rd, "{f \otimes f}"'] \arrow[r, "m"] & L \arrow[d, "f"] \\
		& E.
	\end{tikzcd}\]
	Terlihat bahwa $\tilde a_{ff}=\identity$. Bersama
	\eqref{eqn:cocycle-condition-descent}, ini menunjukkan bahwa untuk setiap
	$(E,f,g)$ berlaku $\tilde a_{fg}^{-1}=\tilde a_{gf}$. Khususnya, ``versi
	universal'' $\tilde a$ juga merupakan isomorfisme. Selesailah bukti.
\end{proof}

% segment-id: o014.li2.chapter7.module-descent.th002
\begin{theorem}[Bentuk kedua penurunan datar]\label{prop:flat-descent-2}
	Misalkan $L$ aljabar-$K$ komutatif yang datar setia. Maka
	% segment-id: o014.li2.chapter7.module-descent.q017
	\[ L \dotimes{K} (\cdot): K\dcate{Mod} \to \mathrm{Desc}_{K \to L} \]
	merupakan ekuivalensi kategori. Salah satu kuasi-inversnya,
	$\tilde{\mathbb{L}}$, dapat diambil sebagai penyama
	% segment-id: o014.li2.chapter7.module-descent.q018
	\[\begin{tikzcd}
		0 \arrow[r] & \tilde{\mathbb{L}}(M, \tilde{a}) \arrow[r] & M \arrow[shift left, r, "a"] \arrow[shift right, r, "{m \mapsto 1 \otimes m}"'] & L \dotimes{K} M,
	\end{tikzcd}\]
	dengan $a$ pemetaan yang bersesuaian dengan $\tilde a$ menurut
	\eqref{eqn:flat-descent-map-aux}.
\end{theorem}
% segment-id: o014.li2.chapter7.module-descent.pf003
\begin{proof}
	Berdasarkan Lema~\ref{prop:descent-data-vs-comodule}, pernyataan ini hanya
	merupakan perumusan ulang Teorema~\ref{prop:flat-descent} tentang penurunan
	datar.
\end{proof}

% segment-id: o014.li2.chapter7.module-descent.p013
Teorema~\ref{prop:flat-descent-2} adalah rumusan penurunan datar yang lazim
dalam geometri aljabar. Di sana, homomorfisme datar setia $K\to L$
bersesuaian dengan suatu peliputan objek geometri dalam arti tertentu.

% segment-id: o014.li2.chapter7.module-descent.p014
Data $\tilde a$ atau $a$ yang bersesuaian kadang-kadang juga dinyatakan
sebagai homomorfisme modul-$L\dotimes{K}L$ kiri, yaitu homomorfisme
bimodul-$(L,L)$
% segment-id: o014.li2.chapter7.module-descent.q019
\begin{align*}
	\overline{a}: M \dotimes{K} L & \to L \dotimes{K} M \\
	m \otimes \ell & \mapsto \sum_{i=1}^n \ell_i \otimes \ell m_i,
\end{align*}
dengan tetap mengandaikan
$a(m)=\sum_{i=1}^n\ell_i\otimes m_i$. Bentuk ini diperoleh dengan
mengontraksikan hasil kali tensor dalam \eqref{eqn:flat-descent-map-aux} secara
tepat dan memberikan bentuk yang lebih ringkas. Versi terkontraksi untuk
$\tilde a_{ji}$ diperoleh dengan cara serupa.

% segment-id: o014.li2.chapter7.module-descent.r002
\begin{remark}[Penurunan datar bagi aljabar dan modul]\label{rem:flat-descent-alg}
	Kita dapat membahas lebih lanjut bagaimana aljabar-$L$, atau modul di atas
	aljabar tersebut, diturunkan ke $K$. Secara konkret, misalkan
	$(M,a),(M',a')\in\Obj(L\dcate{Mod}^{\mathbf{L}})$ adalah citra dari
	$N,N'\in\Obj(K\dcate{Mod})$. Ingat bahwa
	$L\dotimes{K}(\mathord\cdot):K\dcate{Mod}\to L\dcate{Mod}$ adalah funktor
	monoidal \cite[Proposisi 6.6.10]{Li1}; yakni terdapat isomorfisme kanonik
	seperti
	% segment-id: o014.li2.chapter7.module-descent.q020
	\begin{equation}\label{eqn:change-of-ring-monoidal}\begin{aligned}
			(L \dotimes{K} N_1) \dotimes{L} (L \dotimes{K} N_2) & \rightiso L \dotimes{K} (N_1 \dotimes{K} N_2) \\
			(\ell_1 \otimes x) \otimes (\ell_2 \otimes y) & \mapsto \ell_1 \ell_2 \otimes (x \otimes y).
	\end{aligned}\end{equation}
	Gunakan ini untuk memberikan $M\dotimes{L}M'$ struktur komodul-
	$\mathbf{L}$ yang berasal dari $N\dotimes{K}N'$. Sekarang ambil $N'=N$.
	Memberikan struktur aljabar-$K$ pada $N$ sama artinya dengan memberikan
	struktur aljabar-$L$ pada $M$ sedemikian sehingga perkalian
	$M\dotimes{L}M\to M$ dan satuan $L\to M$ keduanya merupakan morfisme dalam
	$L\dcate{Mod}^{\mathbf{L}}$; semua sifat seperti asosiativitas yang
	diperlukan bagi aljabar-$K$ berasal dari $M$. Penurunan modul-$M$ ke
	modul-$N$ harus dipandang dengan cara yang sama. Untuk tujuan ini, data
	penurunan berbentuk \eqref{eqn:flat-descent-map-aux} mungkin lebih praktis
	daripada komodul bagi $\mathbf{L}$, karena funktor yang terlibat semuanya
	monoidal pada $M$, sedangkan $\mathbf{L}$ sendiri bukan funktor monoidal.
\end{remark}
